% =====================================================================
% Physics Notation Accessibility Test Document
% =====================================================================
% Purpose: Test how MathJax 4 (and downstream screen readers like
% NVDA+MathCAT) handle common notation from undergraduate physics
% (first 2 years: classical mechanics, E&M, thermodynamics, waves,
% intro modern physics/QM, intro relativity).
%
% Compile with: pdflatex physics-accessibility-test.tex
%
% For Pandoc conversion:
%   pandoc physics-accessibility-test.tex --mathjax -s \
%       -o physics-accessibility-test.html
%
% IMPORTANT: For the HTML output, MathJax must be configured to load
% the `physics` package extension. Add to your MathJax config:
%   loader: {load: ['[tex]/physics']},
%   tex: {packages: {'[+]': ['physics']}}
%
% Each test case has a comment describing what is being tested and
% the expected speech behavior with MathCAT (verbose mode for clarity).
% =====================================================================

\documentclass[11pt]{article}

\usepackage[utf8]{inputenc}
\usepackage[T1]{fontenc}
\usepackage{amsmath, amssymb, amsthm}
\usepackage{physics}    % key package for physics notation
\usepackage{siunitx}    % SI units
\usepackage{bm}         % bold math
\usepackage{geometry}
\geometry{margin=1in}

\title{Physics Notation Accessibility Test \\
       \large Undergraduate Years 1--2}
\author{ATC Accessibility Pipeline}
\date{}

\begin{document}
\maketitle

\section*{How to use this document}

Each numbered item tests a specific notation pattern. The visual
rendering (in PDF) is the ground truth for what the equation looks like.
For the accessibility test, convert this document to HTML via Pandoc with
MathJax, then read each equation with your screen reader (NVDA+MathCAT
recommended). Note discrepancies between what is seen and what is heard.

% =====================================================================
\section{Classical Mechanics}
% =====================================================================

\begin{enumerate}

\item \textbf{Newton's second law (vector form).}
Tests bold vector notation and the equality of force and rate of change
of momentum.
\[
  \vb{F} = m\vb{a} = \dv{\vb{p}}{t}
\]
% Expected speech: "F equals m a equals the derivative of p with respect to t"
% Watch for: whether bold vectors are announced as "vector F" or just "F";
% whether the derivative is read naturally or as "d p over d t".

\item \textbf{Newton's dot notation for time derivatives.}
Velocity and acceleration as first and second time derivatives. This is
a major source of accessibility issues because the dot has multiple
possible meanings.
\[
  \vb{v} = \dot{\vb{r}}, \qquad
  \vb{a} = \ddot{\vb{r}} = \dv[2]{\vb{r}}{t}
\]
% Expected: "v equals r dot ... a equals r double dot equals second
% derivative of r with respect to t"
% Watch for: does MathCAT say "the time derivative of r" or just "r dot"?

\item \textbf{Work as a line integral.}
Tests the dot product and line integral notation along a path.
\[
  W = \int_C \vb{F} \vdot \dd{\vb{r}}
\]
% Expected: "W equals the integral over C of F dot d r"

\item \textbf{Kinetic energy and angular momentum.}
Tests fractions, exponents, and the cross product.
\[
  T = \frac{1}{2}m\abs{\vb{v}}^2, \qquad
  \vb{L} = \vb{r} \cross \vb{p}
\]
% Expected: "T equals one half m absolute value of v squared ...
% L equals r cross p"

\item \textbf{Rotational dynamics.}
Tests Greek letters (tau, omega, alpha) in a physics context.
\[
  \vb{\tau} = \vb{r} \cross \vb{F} = I\vb{\alpha},
  \qquad
  \vb{L} = I\vb{\omega}
\]
% Watch for: are Greek letters announced clearly? Does MathCAT distinguish
% the role of tau as torque from tau used elsewhere (e.g., proper time)?

\item \textbf{Simple harmonic motion.}
Tests phase, trigonometric functions, and time-dependent expressions.
\[
  x(t) = A\cos(\omega t + \varphi)
\]
% Expected: "x of t equals A cosine of omega t plus phi"
% Watch for: is "phi" announced as "phi" or "varphi"? (\varphi vs \phi)

\item \textbf{Lagrangian and Euler-Lagrange equation.}
Tests calligraphic letters (common for Lagrangian) and partial derivatives.
\[
  \mathcal{L} = T - V, \qquad
  \dv{t}\left(\pdv{\mathcal{L}}{\dot{q}_i}\right)
  - \pdv{\mathcal{L}}{q_i} = 0
\]
% Watch for: is \mathcal{L} announced as "script L" or just "L"?
% Does the partial-with-respect-to-q-dot get read correctly?

\item \textbf{Centripetal acceleration.}
Tests inline math with subscripts and fractions.
The centripetal acceleration is $a_c = v^2/r = \omega^2 r$, directed
toward the center of the circle.
% Watch for: subscript "c" announced clearly; v squared over r vs.
% v over r squared (precedence matters).

\end{enumerate}

% =====================================================================
\section{Electromagnetism}
% =====================================================================

\begin{enumerate}

\item \textbf{Coulomb's law.}
Tests vector notation, hat notation for unit vectors, and exponents.
\[
  \vb{F}_{12} = \frac{1}{4\pi\varepsilon_0}\,
  \frac{q_1 q_2}{r^2}\,\vu{r}_{12}
\]
% Watch for: subscript "12" read as "one two" or "twelve"; unit vector
% r-hat announced clearly; epsilon-naught read naturally.

\item \textbf{Electric field from a potential.}
Tests the gradient operator from the physics package.
\[
  \vb{E} = -\grad V
\]
% Expected: "E equals minus gradient of V" (with the physics package,
% \grad produces a recognizable gradient operator).

\item \textbf{Gauss's law (integral form).}
Tests closed surface integral notation and the dot product.
\[
  \oint_S \vb{E} \vdot \dd{\vb{A}} = \frac{Q_{\text{enc}}}{\varepsilon_0}
\]
% Watch for: is the closed integral (\oint) announced distinctly from
% a regular integral? "Surface integral of E dot d A"?

\item \textbf{Maxwell's equations (differential form).}
Tests divergence, curl, and partial time derivatives together.
\begin{align*}
  \div \vb{E} &= \frac{\rho}{\varepsilon_0}
  & \div \vb{B} &= 0 \\
  \curl \vb{E} &= -\pdv{\vb{B}}{t}
  & \curl \vb{B} &= \mu_0 \vb{J}
    + \mu_0 \varepsilon_0 \pdv{\vb{E}}{t}
\end{align*}
% Watch for: aligned equations; "del dot E" vs "divergence of E";
% navigation between the four equations.

\item \textbf{Lorentz force law.}
Tests addition of two vector terms with a cross product.
\[
  \vb{F} = q(\vb{E} + \vb{v} \cross \vb{B})
\]
% Expected: "F equals q times the quantity E plus v cross B"

\item \textbf{Wave equation for EM waves in vacuum.}
Tests the Laplacian and second time derivatives.
\[
  \laplacian \vb{E} = \mu_0 \varepsilon_0 \pdv[2]{\vb{E}}{t}
\]
% Watch for: \laplacian announced as "Laplacian of E" or as
% "del squared E"; second partial derivative read correctly.

\item \textbf{Poynting vector.}
Tests cross product in the context of energy flux.
\[
  \vb{S} = \frac{1}{\mu_0}\,\vb{E} \cross \vb{B}
\]

\item \textbf{Capacitor energy.}
Tests inline math with a fraction.
The energy stored in a capacitor is $U = \frac{1}{2}CV^2$, where $C$ is
capacitance and $V$ is the voltage across the plates.

\end{enumerate}

% =====================================================================
\section{Thermodynamics and Statistical Mechanics}
% =====================================================================

\begin{enumerate}

\item \textbf{Ideal gas law.}
Tests subscripts, multiplication, and standard variable names.
\[
  PV = nRT = Nk_B T
\]
% Watch for: is "k_B" announced as "k sub B" or "Boltzmann constant"?
% MathCAT may or may not have a rule for this.

\item \textbf{First law of thermodynamics.}
Tests inexact differentials (often written with $\d$ vs $\delta$ or with
$\dbar$ in some conventions).
\[
  \dd{U} = \delta Q - \delta W
\]
% Watch for: distinction between "d U" (exact) and "delta Q" (inexact)?
% Most speech engines do not distinguish these visually.

\item \textbf{Thermodynamic identity.}
Tests multiple differentials in one expression.
\[
  \dd{U} = T\dd{S} - P\dd{V} + \mu\dd{N}
\]

\item \textbf{Entropy (Boltzmann).}
Tests the natural log and the Greek letter Omega for microstates.
\[
  S = k_B \ln \Omega
\]
% Watch for: is "ln" announced as "natural log" or "L N"?

\item \textbf{Partition function and free energy.}
Tests summation notation, exponentials, and the inverse temperature.
\[
  Z = \sum_i e^{-\beta E_i}, \qquad
  F = -k_B T \ln Z
\]
% Expected: "Z equals the sum over i of e to the minus beta E sub i"

\item \textbf{Maxwell-Boltzmann distribution.}
Tests a complex expression with multiple operations.
\[
  f(v) = 4\pi n
  \left(\frac{m}{2\pi k_B T}\right)^{3/2}
  v^2 e^{-mv^2/(2k_B T)}
\]
% Watch for: fractional exponent 3/2; nested fraction inside the exponent.

\end{enumerate}

% =====================================================================
\section{Waves and Oscillations}
% =====================================================================

\begin{enumerate}

\item \textbf{Wave equation in one dimension.}
Tests second partial derivatives in space and time.
\[
  \pdv[2]{y}{t} = v^2 \pdv[2]{y}{x}
\]
% Watch for: are both second partials announced consistently?

\item \textbf{Plane wave solution.}
Tests trig with combined arguments.
\[
  y(x,t) = A\cos(kx - \omega t + \varphi)
\]
% Watch for: variable order and signs in argument; phase phi.

\item \textbf{Dispersion relation and group velocity.}
Tests derivatives of one variable with respect to another.
\[
  v_p = \frac{\omega}{k}, \qquad
  v_g = \dv{\omega}{k}
\]
% Watch for: subscripts p and g; whether $v_g$ is identified as
% "group velocity" or just "v sub g".

\item \textbf{Standing wave.}
Tests product of two trig functions.
\[
  y(x,t) = 2A \sin(kx) \cos(\omega t)
\]

\item \textbf{Doppler effect.}
Tests fractions with multiple terms and signs.
\[
  f_{\text{obs}} = f_{\text{src}}
  \left(\frac{v \pm v_{\text{obs}}}{v \mp v_{\text{src}}}\right)
\]
% Watch for: $\pm$ and $\mp$ announced clearly; text subscripts.

\end{enumerate}

% =====================================================================
\section{Introduction to Quantum Mechanics}
% =====================================================================

\begin{enumerate}

\item \textbf{de Broglie wavelength.}
Tests Greek letters and a simple ratio.
\[
  \lambda = \frac{h}{p}
\]

\item \textbf{Schr\"{o}dinger equation (time-dependent).}
Tests $i$, $\hbar$, operator hats, and partial derivatives all together.
\[
  i\hbar\,\pdv{\Psi}{t} = \hat{H}\Psi
\]
% Watch for: i (imaginary unit) read as "i" not "1"; hbar as "h-bar";
% \hat{H} as "H hat" or "Hamiltonian operator".

\item \textbf{Time-independent Schr\"odinger equation.}
\[
  \hat{H}\psi = E\psi, \qquad
  \hat{H} = -\frac{\hbar^2}{2m}\laplacian + V(\vb{r})
\]
% Watch for: distinction between uppercase Psi and lowercase psi.

\item \textbf{Position-space momentum operator.}
Tests the gradient operator combined with i and hbar.
\[
  \hat{\vb{p}} = -i\hbar\,\grad
\]
% Watch for: is "p-hat" or "momentum operator" announced? bold p-hat
% is a particularly compound notation.

\item \textbf{Expectation values (bra-ket).}
Tests bra-ket notation from the physics package.
\[
  \ev{\hat{A}} = \mel{\psi}{\hat{A}}{\psi}
  = \int \psi^* \hat{A} \psi \dd{x}
\]
% Expected (MathCAT): "expectation value of A hat equals
% matrix element bra psi A hat ket psi equals..."
% Watch for: is the bra-ket notation recognized? Does MathCAT say
% "the bra psi", "ket psi", "matrix element"?

\item \textbf{Probability density.}
Tests complex conjugate (the asterisk) and absolute value.
\[
  P(x) = \abs{\psi(x)}^2 = \psi^*(x)\psi(x)
\]
% Watch for: $\psi^*$ as "psi star" or "psi conjugate"?

\item \textbf{Commutators.}
Tests the commutator bracket notation (a critical QM idiom).
\[
  \comm{\hat{x}}{\hat{p}} = i\hbar, \qquad
  \comm{\hat{L}_i}{\hat{L}_j} = i\hbar\,\epsilon_{ijk}\hat{L}_k
\]
% Watch for: does MathCAT announce "the commutator of x and p" or
% does it read "open bracket x hat comma p hat close bracket"?
% The second equation also tests Levi-Civita epsilon with three indices.

\item \textbf{Uncertainty principle.}
Tests the capital delta (uncertainty) and an inequality.
\[
  \Delta x \, \Delta p \geq \frac{\hbar}{2}
\]

\item \textbf{Particle in a box: energy eigenvalues.}
Tests subscripted energy levels.
\[
  E_n = \frac{n^2 \pi^2 \hbar^2}{2mL^2}, \qquad n = 1, 2, 3, \ldots
\]

\item \textbf{Particle in a box: wavefunctions.}
Tests square roots, sines, and the use of $n$ as both index and parameter.
\[
  \psi_n(x) = \sqrt{\frac{2}{L}}\,\sin\!\left(\frac{n\pi x}{L}\right)
\]

\end{enumerate}

% =====================================================================
\section{Introduction to Special Relativity}
% =====================================================================

\begin{enumerate}

\item \textbf{Lorentz factor.}
Tests nested radical and fraction.
\[
  \gamma = \frac{1}{\sqrt{1 - v^2/c^2}}
\]
% Watch for: nested fraction within a radical; precedence of operations.

\item \textbf{Time dilation and length contraction.}
Tests primed and unprimed variables (a common confusion).
\[
  \Delta t = \gamma \Delta t', \qquad
  L = L_0/\gamma
\]
% Watch for: is the prime distinguished in speech? "delta t prime"
% versus "delta t"?

\item \textbf{Mass-energy and energy-momentum relation.}
Tests the famous equations together.
\[
  E = mc^2, \qquad
  E^2 = (pc)^2 + (mc^2)^2
\]

\item \textbf{Four-vector notation (intro).}
Tests indexed four-vectors with Greek index conventions.
\[
  x^\mu = (ct, x, y, z), \qquad
  p^\mu = (E/c, \vb{p})
\]
% Watch for: indexed notation; the upper Greek index mu read clearly.

\item \textbf{Spacetime interval.}
Tests the Minkowski metric (sums of squares with one negative sign).
\[
  \dd{s}^2 = -c^2 \dd{t}^2 + \dd{x}^2 + \dd{y}^2 + \dd{z}^2
\]
% Note: sign convention varies; this is the "mostly plus" convention.
% Watch for: are the differentials read consistently?

\end{enumerate}

% =====================================================================
\section{Atomic and Nuclear Physics (Introductory)}
% =====================================================================

\begin{enumerate}

\item \textbf{Bohr model energy levels.}
Tests subscript, negative number, and units.
\[
  E_n = -\frac{13.6\;\text{eV}}{n^2}
\]
% Watch for: how is the unit "eV" announced? "e V" or "electron-volts"?

\item \textbf{Rydberg formula.}
Tests two terms with different subscripts in a difference.
\[
  \frac{1}{\lambda} = R_H \left(\frac{1}{n_1^2} - \frac{1}{n_2^2}\right)
\]

\item \textbf{Radioactive decay law.}
Tests exponential decay.
\[
  N(t) = N_0\, e^{-\lambda t}
\]

\item \textbf{Half-life relation.}
Tests natural log and the structure of a defining equation.
\[
  t_{1/2} = \frac{\ln 2}{\lambda}
\]
% Watch for: subscript 1/2 announced as "one-half" or "1 slash 2"?

\end{enumerate}

% =====================================================================
\section{SI Units (via siunitx)}
% =====================================================================

These test how units are announced when typeset with the \texttt{siunitx}
package. Note: Pandoc may or may not preserve \texttt{siunitx} commands;
MathJax has partial support for \texttt{siunitx} via an extension.

\begin{enumerate}

\item A speed: \SI{3.00e8}{\meter\per\second}.
% Watch for: is "m/s" announced as "meters per second"? Is the exponential
% notation handled?

\item An energy: \SI{1.602e-19}{\joule}.

\item An angle: \SI{45}{\degree}.

\item A combined unit: a magnetic flux density of
\SI{1.5}{\tesla} or equivalently \SI{1.5}{\kilogram\per\ampere\per\second\squared}.

\item A range: \SIrange{0}{100}{\celsius}.

\end{enumerate}

% =====================================================================
\section{Mixed Notation Stress Tests}
% =====================================================================

These combine multiple notation challenges in a single expression to
test how the speech engine handles compound complexity.

\begin{enumerate}

\item \textbf{Aligned equation with continuation.}
Tests \texttt{align} environment and phantoms for alignment.
\begin{align}
  \vb{F}_{\text{net}}
    &= m\vb{a} \\
  &= m\dv[2]{\vb{r}}{t} \\
  &= -\grad U(\vb{r})
\end{align}
% Watch for: numbered equations; the continuation pattern read as a
% single chain.

\item \textbf{Piecewise function.}
Tests cases environment.
\[
  V(x) = \begin{cases}
    0 & 0 < x < L \\
    \infty & \text{otherwise}
  \end{cases}
\]
% Watch for: structural reading of cases.

\item \textbf{Matrix (Pauli sigma-z).}
Tests 2x2 matrices.
\[
  \sigma_z = \begin{pmatrix} 1 & 0 \\ 0 & -1 \end{pmatrix}
\]
% Watch for: matrix dimensions announced; navigation by rows/columns.

\item \textbf{Definite integral with limits and complicated integrand.}
\[
  \langle x \rangle =
  \int_{-\infty}^{\infty} \psi^*(x)\, x\, \psi(x)\, \dd{x}
\]
% Watch for: infinite limits; product of three things being integrated.

\item \textbf{A typical homework expression.}
Tests realistic compound notation.
\[
  \vb{B}(\vb{r}) = \frac{\mu_0}{4\pi} \int_C
  \frac{I\,\dd{\vb{r}'} \cross (\vb{r} - \vb{r}')}{\abs{\vb{r} - \vb{r}'}^3}
\]
% Biot-Savart law. Watch for: primed and unprimed positions distinguished;
% absolute value with cube exponent; cross product inside the integrand.

\end{enumerate}

\end{document}